% Generated draft: requires mathematical and source review.
\documentclass[11pt]{article}
\usepackage[margin=25mm]{geometry}
\usepackage{fontspec}
\setmainfont{Latin Modern Roman}
\usepackage{amsmath,amssymb,mathtools,unicode-math}
\setmathfont{Latin Modern Math}
\usepackage{enumitem,needspace,xurl,hyperref}
\setlength{\parindent}{0pt}
\setlength{\parskip}{5pt}
\setlength{\emergencystretch}{3em}
\newcommand{\N}{\mathbb{N}}
\newcommand{\Z}{\mathbb{Z}}
\newcommand{\Q}{\mathbb{Q}}
\newcommand{\R}{\mathbb{R}}
\newcommand{\C}{\mathbb{C}}
\newcommand{\F}{\mathbb{F}}
\newcommand{\A}{\mathbb{A}}
\newcommand{\PP}{\mathbb{P}}
\newcommand{\E}{\mathbb{E}}
\newcommand{\eps}{\varepsilon}
\newcommand{\dd}{\,\mathrm{d}}
\begin{document}
\section*{UnsolvedMath ID: 30002582}
\textbf{Source problem code:} catalogue entry
\subsection{Definitions and mathematical statement}
The supplied record poses several related questions. Its terminology is not sufficiently specified to give all of them unconditional precise formulations. The following notation uses the usual Drury--Arveson framework for multiplier varieties; this framework is a contextual convention requiring confirmation against the cited report.

For \(d\in\N\), let \(\mathbb B_d=\{z\in\C^d:\|z\|<1\}\), and let \(\mathbb B_\infty\) be the open unit ball of \(\ell^2(\N)\). Write \(\mathbb D=\{\zeta\in\C:|\zeta|<1\}\). For \(d\in\N\cup\{\infty\}\), let \(H_d^2\) be the reproducing kernel Hilbert space on \(\mathbb B_d\) with kernel
\[
 k_d(z,w)=\frac{1}{1-\langle z,w\rangle},
 \qquad \langle z,w\rangle=\sum_j z_j\overline{w_j}.
\]
For a subset \(V\subseteq\mathbb B_d\), let \(H_V\) be the reproducing kernel Hilbert space with kernel \(k_d|_{V\times V}\), and define
\[
 \mathcal M_V=\operatorname{Mult}(H_V)
 =\{\varphi:V\to\C: M_\varphi:H_V\to H_V\text{ is bounded}\},
 \qquad (M_\varphi h)(v)=\varphi(v)h(v).
\]
Equip this algebra with the multiplier norm \(\|\varphi\|_{\mathcal M_V}=\|M_\varphi\|\), and write \(\mathcal M_d=\operatorname{Mult}(H_d^2)\). In this framework a multiplier variety is a set of the form
\[
 V=Z(\mathcal S)
 :=\{z\in\mathbb B_d:\varphi(z)=0\text{ for every }\varphi\in\mathcal S\},
 \qquad \mathcal S\subseteq\mathcal M_d.
\]
The precise zero-set convention used by the report is not supplied and must be confirmed.

\begin{enumerate}
\item For multiplier varieties, does there exist an intrinsic local characterization of the global zero-set notion above? Are their classical irreducible components themselves multiplier varieties, or varieties in the intended local sense? Here classical irreducible components are understood provisionally as maximal irreducible complex-analytic components, where irreducibility means that the component cannot be written as the union of two proper closed analytic subsets. The record does not specify the local characterization sought, the meaning of intrinsic, the allowable ambient dimensions, or the relation between the proposed local notion and the global multiplier zero-set notion. These omissions prevent a definitive quantified formulation of this part.

\item Let \(V\) and \(W\) be irreducible multiplier varieties in a finite-dimensional ball. If \(V\) and \(W\) are multiplier-biholomorphic, must their multiplier algebras satisfy
\[
 \mathcal M_V\cong\mathcal M_W?
\]
The record does not define multiplier biholomorphism, including its requirements on coordinate multipliers or ambient extensions. It also does not specify whether the algebra isomorphism sought is algebraic, bounded with bounded inverse, completely bounded with completely bounded inverse, or completely isometric. The ambient dimensions of the two varieties and the intended notion of irreducibility require confirmation. The displayed implication is therefore conditional on these missing conventions.

\item Is multiplier biholomorphism an equivalence relation on the intended class of multiplier varieties? Explicitly, for the intended relation \(\sim_{\mathrm{mb}}\), are
\[
 V\sim_{\mathrm{mb}}V,\qquad
 V\sim_{\mathrm{mb}}W\Longrightarrow W\sim_{\mathrm{mb}}V,
\]
\[
 V\sim_{\mathrm{mb}}W\ \text{and}\ W\sim_{\mathrm{mb}}U
 \Longrightarrow V\sim_{\mathrm{mb}}U
\]
valid for all varieties in that class? Neither the exact relation nor the class, including any dimension restrictions, is specified by the supplied record.

\item Let \(f:\mathbb D\to\mathbb B_\infty\) be a proper embedding with \(f(0)=0\), and put \(V=f(\mathbb D)\). Interpreting the embedding as holomorphic and the character algebra as \(\mathcal M_V\), let \(z_j:V\to\C\) denote the restricted coordinate multiplier \(z_j(v)=v_j\). A character means a nonzero complex-linear multiplicative functional \(\chi:\mathcal M_V\to\C\). Define evaluation at the image origin by \(\rho_0(\varphi)=\varphi(0)\). Is
\[
 \bigl\{\chi:\chi\text{ is a character of }\mathcal M_V,
       \ \chi(z_j)=0\text{ for every }j\in\N\bigr\}
 =\{\rho_0\}?
\]
This formulation retains all the displayed hypotheses on \(f\), but its interpretation is conditional: the record does not explicitly identify the character algebra or define proper embedding. In the infinite-dimensional ball one must distinguish topological properness, meaning compact inverse images of compact sets, from the boundary condition \(\|f(\zeta)\|\to1\) as \(|\zeta|\to1\); the supplied record does not select or relate these conditions. It also does not state separately that \(f(\mathbb D)\) is a multiplier variety.
\end{enumerate}

These are statement-only questions. The supplied open-status label and publication-verification label do not establish their current mathematical status or resolve the formulation issues above.
\subsection{Short English statement}
The record asks whether multiplier varieties admit an intrinsic local description, whether classical irreducible components remain varieties, whether multiplier biholomorphism of irreducible varieties in finite-dimensional balls forces an isomorphism of multiplier algebras, and whether multiplier biholomorphism is an equivalence relation. It also asks whether evaluation at the origin is the unique character annihilating every coordinate multiplier for the image of a proper disk embedding into \(\mathbb B_\infty\) fixing the origin. Several definitions and hypotheses require source confirmation.
\subsection{Sources}
\begin{itemize}
\item Supplied local UnsolvedMath registry record, permanent ID \(30002582\), title \emph{Local Geometry and Algebraic Classification of Multiplier Varieties}, source problem code \texttt{catalogue entry}, collection \emph{Oberwolfach Reports Open Problems}. No individual proposer is supplied. The proposed-year field is \(2014\).
\item Supplied report citation: \emph{Hilbert Modules and Complex Geometry (2015)}, Oberwolfach Reports, pp.~1139--1219, DOI \texttt{10.4171/owr/2014/21}: \url{https://doi.org/10.4171/owr/2014/21}. The background describes Oberwolfach Reports as documenting mathematical workshops and their open-problem discussions. No problem-specific page or speaker attribution is supplied. The citation's year \(2015\), the proposed year \(2014\), and the DOI containing \(2014\) are retained as supplied.
\item Supplied literature reference: \emph{Operator algebras for analytic varieties}, \url{https://arxiv.org/abs/1201.4072}.
\item Supplied literature reference: \emph{Complete-Pick isomorphism problem}, \url{https://arxiv.org/abs/1211.1116}.
\item The supplied literature-source list also includes the Oberwolfach report: \url{https://doi.org/10.4171/owr/2014/21}.
\item Literature assessment dated exactly \texttt{2026-08-21}: the supplied curation labels the problem open and states that isomorphism implies multiplier biholomorphism in broad finite-dimensional classes, that homogeneous cases are understood, and that the listed local-variety, converse-isomorphism, equivalence-relation, and character questions retain unresolved general cases. This is supplied dated curation, not an independently verified assessment. The background explicitly requests independent verification of time-sensitive status and cited claims.
\item The registry's statement-status label is \texttt{corrected\_verified}; its accompanying explanation says mathematical content was preserved while notation, typography, or omitted definitions were repaired for publication. This supplied editorial label does not resolve the missing definitions or constitute independent mathematical review of this draft.
\end{itemize}
\end{document}
